% generated by 07_manuscript/make_numbers.py from conclusions.md and DEVIATIONS.md (verbatim; LaTeX escaping only)
\newcommand{\concHK}{Under free-living wear (WildPPG, 13 held-out persons), 8.80 of 39 (person, state) cells exceeded 2 false alarms per monitored hour (seed average), more than the 4.39 expected if the per-cell tolerance guarantee held (person-cluster bootstrap lower bound > 0). The guarantee of Proposition 2a did not transfer to these data. Its assumption A(0) is violated in at least part of the cells. \S4.2 states which source of the violation the data support. All participants of this dataset followed the same one-day route with an altitude change of about 3,000 m, so the persons are not independent in the direction of any change over the day, and the result may not carry over to free-living wear without such a change.}
\newcommand{\concHKBranch}{rejected}
\newcommand{\concHKReason}{X2 bootstrap 5 \% quantile 1.03 > 0}
\newcommand{\concHKNoLone}{Under free-living wear (WildPPG, 13 held-out persons), 8.80 of 39 (person, state) cells exceeded 2 false alarms per monitored hour (seed average), more than the 4.39 expected if the per-cell tolerance guarantee held (person-cluster bootstrap lower bound > 0). The guarantee of Proposition 2a did not transfer to these data. Its assumption A(0) is violated in at least part of the cells. \S4.2 states which source of the violation the data support.}
\newcommand{\concHKone}{The excess remains when lag-1 dependence is allowed. All participants of this dataset followed the same one-day route with an altitude change of about 3,000 m, so the persons are not independent in the direction of any change over the day, and the result may not carry over to free-living wear without such a change.}
\newcommand{\concHKoneBranch}{rejected}
\newcommand{\concHKoneReason}{H-K rejected; H-K1 quantile 6.54 > 0}
\newcommand{\concHKoneNoLone}{The excess remains when lag-1 dependence is allowed.}
\newcommand{\concAttr}{Block scores are not exchangeable in time order, and exceedances were associated with altitude change since calibration (non-stationarity over the day); physiological and behavioural components were not separated. All participants of this dataset followed the same one-day route with an altitude change of about 3,000 m, so the persons are not independent in the direction of any change over the day, and the result may not carry over to free-living wear without such a change.}
\newcommand{\concAttrBranch}{row 3}
\newcommand{\concAttrReason}{p\_order 0.004 < 0.05; dAlt 90 \% CI [0.14, 0.88] above 0}
\newcommand{\concAttrNoLone}{Block scores are not exchangeable in time order, and exceedances were associated with altitude change since calibration (non-stationarity over the day); physiological and behavioural components were not separated.}
\newcommand{\concHR}{Rolling calibration kept the long-run rate consistent with its expectation bound of 0.857/h (mean of persons 0.88, 90 \% CI 0.74–1.03): not refuted at n = 13, which is not the same as ``the bound holds''. Rolling calibration gives only an expectation bound, not the per-cell tolerance guarantee. All participants of this dataset followed the same one-day route with an altitude change of about 3,000 m, so the persons are not independent in the direction of any change over the day, and the result may not carry over to free-living wear without such a change.}
\newcommand{\concHRBranch}{not rejected}
\newcommand{\concHRReason}{t p 0.383 $\geq$ 0.05}
\newcommand{\concHRNoLone}{Rolling calibration kept the long-run rate consistent with its expectation bound of 0.857/h (mean of persons 0.88, 90 \% CI 0.74–1.03): not refuted at n = 13, which is not the same as ``the bound holds''. Rolling calibration gives only an expectation bound, not the per-cell tolerance guarantee.}
\newcommand{\concRRthin}{The paired comparison does not separate short-range dependence from non-stationarity. All participants of this dataset followed the same one-day route with an altitude change of about 3,000 m, so the persons are not independent in the direction of any change over the day, and the result may not carry over to free-living wear without such a change.}
\newcommand{\concRRthinBranch}{CI contains 0}
\newcommand{\concRRthinReason}{90 \% CI [-0.22, 0.0] contains 0}
\newcommand{\concRRthinNoLone}{The paired comparison does not separate short-range dependence from non-stationarity.}
\newcommand{\sentLone}{All participants of this dataset followed the same one-day route with an altitude change of about 3,000 m, so the persons are not independent in the direction of any change over the day, and the result may not carry over to free-living wear without such a change.}
\newcommand{\sentDtwo}{The model intended to separate altitude change from elapsed time did not meet the pre-specified stability rule (14.4 \% of bootstrap draws dropped; limit 5 \%). Altitude change and elapsed time since calibration were therefore not separated in these data. This case was not covered by the analysis plan; the sentence was added after unblinding.}
\newcommand{\sentFtwo}{Rolling calibration kept the long-run rate consistent with its expectation bound of 0.857/h (mean of persons 0.88, 90 \% CI 0.74–1.03): not refuted at n = 13, which is not the same as ``the bound holds''. This came at a sensitivity cost of 0.243 / 0.373 / 0.453 vs 0.371 / 0.475 / 0.502 for K (+20 / +30 / +40 bpm) and with masking of a second event 2–7 state-blocks after a first one (detection 0.192–0.261 vs 0.272–0.391; none at h = 1 and h $\geq$ 8). Rolling calibration gives only an expectation bound, not the per-cell tolerance guarantee.}
